% markpage — export LaTeX (cible : xelatex ou lualatex).
% La compilation avec pdflatex échouera : fontspec + UTF-8 natif
% dans les blocs de code requièrent xelatex / lualatex.
\documentclass[11pt,a4paper]{article}
\usepackage{fontspec}
\IfFontExistsTF{DejaVu Serif}{\setmainfont{DejaVu Serif}}{}
\IfFontExistsTF{DejaVu Sans}{\setsansfont{DejaVu Sans}}{}
\IfFontExistsTF{DejaVu Sans Mono}{\setmonofont{DejaVu Sans Mono}}{}

\IfFileExists{french.ldf}{\usepackage[french]{babel}}{}
\usepackage{amsmath,amssymb,amsthm}
\usepackage{stmaryrd}
\usepackage{graphicx}
\usepackage[export]{adjustbox}
\usepackage{hyperref}
\usepackage{xcolor}
\usepackage{listings}
\usepackage{enumitem}
\usepackage{booktabs}
\usepackage{tabularx}
\usepackage[normalem]{ulem}
\usepackage[breakable,skins]{tcolorbox}
\usepackage{newunicodechar}

% DejaVu Serif (our main text font) doesn't ship the Unicode
% math-bracket / logic glyphs that show up in our prose (e.g.
% explanations of the editor's ligatures). Redirect each known
% glyph to its LaTeX command via \ensuremath, so it renders
% from the math font even when typed in normal text.
\newunicodechar{⟦}{\ensuremath{\llbracket}}
\newunicodechar{⟧}{\ensuremath{\rrbracket}}
\newunicodechar{⟨}{\ensuremath{\langle}}
\newunicodechar{⟩}{\ensuremath{\rangle}}
\newunicodechar{⊢}{\ensuremath{\vdash}}
\newunicodechar{⊣}{\ensuremath{\dashv}}
\newunicodechar{⊨}{\ensuremath{\models}}
\newunicodechar{⊥}{\ensuremath{\bot}}
\newunicodechar{⊤}{\ensuremath{\top}}
\newunicodechar{∀}{\ensuremath{\forall}}
\newunicodechar{∃}{\ensuremath{\exists}}
\newunicodechar{∈}{\ensuremath{\in}}
\newunicodechar{∉}{\ensuremath{\notin}}
\newunicodechar{∅}{\ensuremath{\emptyset}}
\newunicodechar{⊂}{\ensuremath{\subset}}
\newunicodechar{⊆}{\ensuremath{\subseteq}}
\newunicodechar{⊃}{\ensuremath{\supset}}
\newunicodechar{⊇}{\ensuremath{\supseteq}}
\newunicodechar{∪}{\ensuremath{\cup}}
\newunicodechar{∩}{\ensuremath{\cap}}
\newunicodechar{∧}{\ensuremath{\land}}
\newunicodechar{∨}{\ensuremath{\lor}}
\newunicodechar{¬}{\ensuremath{\neg}}
\newunicodechar{★}{\ensuremath{\bigstar}}
\newunicodechar{✓}{\ensuremath{\checkmark}}
\newunicodechar{✗}{\ensuremath{\times}}
\newunicodechar{♥}{\ensuremath{\heartsuit}}
\newunicodechar{♦}{\ensuremath{\diamondsuit}}
\newunicodechar{♠}{\ensuremath{\spadesuit}}
\newunicodechar{♣}{\ensuremath{\clubsuit}}
\newunicodechar{⚠}{\textbf{[!]}}

\hypersetup{colorlinks=true, linkcolor=blue!50!black, urlcolor=blue!50!black}
\lstset{
  basicstyle=\ttfamily\small,
  breaklines=true,
  frame=single,
  framesep=4pt,
}

% Theorem-like environments matching the ::: admonitions (§17 of
% the markpage SPEC). The shared counter `theorem` keeps lemma /
% proposition / corollary numbered in the same sequence as the
% theorems themselves, which is the usual mathematical convention.
\newtheorem{theorem}{Théorème}
\newtheorem{lemma}[theorem]{Lemme}
\newtheorem{proposition}[theorem]{Proposition}
\newtheorem{corollary}[theorem]{Corollaire}
\theoremstyle{definition}
\newtheorem{definition}{Définition}
\newtheorem{example}{Exemple}
\newtheorem{remark}{Remarque}

\title{catdiagram — diagrammes commutatifs déclaratifs}
\author{Test Author \\ Test Org}
\date{2026-01-01}

\begin{document}
\maketitle

Le fence \texttt{```catdiagram} accepte la syntaxe documentée dans CD-SPEC.md
(signature + équations + flèches universelles). Le parser produit du
Mermaid que le pipeline existant rend en SVG.


\section{Triangle commutatif}

\begin{lstlisting}
objects:  A, B, C
morphisms:
  f : A -> B
  g : B -> C
  h : A -> C
equations:
  h = g . f
\end{lstlisting}


Voir \textbackslash{}ref\{fig:triangle\}.


\section{Carré du produit}

\begin{lstlisting}
objects:  X, A, B, P
morphisms:
  pi1 : P -> A
  pi2 : P -> B
  f   : X -> A
  g   : X -> B
induced:
  u : X -> P  by (f, g)
equations:
  pi1 . u = f
  pi2 . u = g
\end{lstlisting}


\section{Coproduit}

\begin{lstlisting}
objects:  X, A, B, S
morphisms:
  i1 : A -> S
  i2 : B -> S
  f  : A -> X
  g  : B -> X
induced:
  v : S -> X  by (f, g)
equations:
  v . i1 = f
  v . i2 = g
\end{lstlisting}


\section{Carré de naturalité}

\begin{lstlisting}
objects:  F(X), F(Y), G(X), G(Y)
morphisms:
  Fh    : F(X) -> F(Y)
  Gh    : G(X) -> G(Y)
  eta_X : F(X) -> G(X)
  eta_Y : F(Y) -> G(Y)
equations:
  Gh . eta_X = eta_Y . Fh
\end{lstlisting}


\section{Pullback}

\begin{lstlisting}
objects:  X, A, B, C, P
morphisms:
  f  : A -> C
  g  : B -> C
  p1 : P -> A
  p2 : P -> B
  h  : X -> A
  k  : X -> B
induced:
  u : X -> P  by (h, k)
equations:
  f . p1 = g . p2
  f . h  = g . k
  p1 . u = h
  p2 . u = k
\end{lstlisting}


\section{Égaliseur (paire parallèle)}

\begin{lstlisting}
objects:  E, A, B, Z
morphisms:
  f : A -> B
  g : A -> B
  e : E -> A
  h : Z -> A
induced:
  u : Z -> E  by (h)
equations:
  f . e = g . e
  f . h = g . h
  e . u = h
\end{lstlisting}


\section{Fonctorialité}

\begin{lstlisting}
objects:  F(A), F(B), F(C)
morphisms:
  Ff   : F(A) -> F(B)
  Fg   : F(B) -> F(C)
  Fgof : F(A) -> F(C)
equations:
  Fgof = Fg . Ff
\end{lstlisting}


\section{Objet terminal}

\begin{lstlisting}
objects:  T, A
induced:
  t : A -> T  by ()
\end{lstlisting}


\section{Avec modificateurs (mono / epi / iso)}

\begin{lstlisting}
direction: LR
objects:  A, B, C
morphisms:
  i : A -> B (mono)
  p : B -> C (epi)
  q : A -> C (iso)
\end{lstlisting}


\section{Cas d'erreur — équation mal typée}

L'équation \texttt{f = g} ci-dessous est mal typée (f va de A vers B, g va de
B vers C). Le typechecker doit la rejeter avec une diagnostic
positionnel.


\begin{lstlisting}
objects:  A, B, C
morphisms:
  f : A -> B
  g : B -> C
equations:
  f = g
\end{lstlisting}

\end{document}
